\documentclass{article}
\usepackage{amsmath}
\usepackage{listings}
\usepackage{microtype}
\begin{document}

\subsection{Kritische Reflexion}
\label{sec:reflexion}

\subsubsection{Physics Engine (Rust)}
\label{sec:reflexion-rust}
% [PLATZHALTER - durch zuständige Teammitglieder zu ergänzen]

\subsubsection{Digitale Signalverarbeitung (Python-DSP-Pipeline)}
\label{sec:reflexion-dsp}
% [PLATZHALTER - durch zuständige Teammitglieder zu ergänzen]

\subsubsection{Geometrie und Visualisierung}
\label{sec:reflexion-geometrie}
% [PLATZHALTER - durch zuständige Teammitglieder zu ergänzen]

\subsubsection{Cluster-Infrastruktur und Parallelisierung}
\label{sec:reflexion-cluster}
Die gewählte Lösung -- statische IP-Adressierung und eine feste,
manuell erstellte Liste von SSH-Aufrufen in \texttt{run\_cluster.sh} --
war für die vorliegende Größenordnung von acht Nodes ausreichend und
pragmatisch umsetzbar, stellt jedoch keinen wartbaren oder skalierbaren
Ansatz dar. Bei einer größeren Anzahl an Nodes oder bei sich ändernden
Node-Adressen wäre eine zentrale Konfigurationsdatei die deutlich
robustere Lösung gewesen. Der erhebliche Zeitaufwand für
Zugangsdaten-Resets und SSH-Fehlkonfiguration deutet zudem darauf hin,
dass die Node-Zugänge zu Beginn des Projekts nicht ausreichend
vorbereitet bzw. dokumentiert waren. Gleichzeitig zeigt sich hier auch
eine Grenze der klassischen Erfolgsmessung über Codezeilen oder
Seitenumfang: Ein wesentlicher Teil des Arbeitsaufwands in diesem Bereich
bestand aus Infrastrukturarbeit, die sich weder im finalen Code noch
unmittelbar im Umfang der Dokumentation niederschlägt, für das
Funktionieren der verteilten Ausführung jedoch notwendig war.

\subsubsection{Reflexion als Team}
\label{sec:reflexion-team}
Über die fachlichen Bereiche hinaus lässt sich auch der organisatorische
Ablauf des Projekts kritisch einordnen. Die iterative, erst im
Projektverlauf entstehende Aufgabenverteilung ermöglichte zwar
Flexibilität, erschwerte jedoch eine verlässliche Zeitplanung von Beginn
an.

\end{document}
